% === Begin Label Data ===
% ID: 2
% 难度星级: 
% 标签: 
% 备注: 
% 组卷引用次数: 0
% === End  Label Data ===

\begin{problem}{2010}{G}{浙江卷（理）}{22}{导数}
已知函数$f(x)=(x-a)^2(x+b)e^x$，$a,b\in\mathbb{R}$，$x=a$是$f(x)$的一个极大值点.

（1）求$b$的取值范围;

（2）设$x_1,x_2,x_3$是$f(x)$的3个极值点，问是否存在实数$b$，可找到$x_4\in\mathbb{R}$，使得$x_1,x_2,x_3,x_4$的某种排列$x_{i_1},x_{i_2},x_{i_3},x_{i_4}$（其中$\{i_1,i_2,i_3,i_4\}=\{1,2,3,4\}$）依次成等差数列?
\end{problem}

\begin{solutions}
(I) $f'(x) = e^x (x-a) \left[x^2+(3-a+b)x+2b-ab-a\right]$,

令 $g(x) = x^2+(3-a+b)x+2b-ab-a$, 则

$\Delta = (3-a+b)^2-4(2b-ab-a) = (a+b-1)^2+8 > 0$, (步骤1)

于是，假设 $x_1, x_2$ 是 $g(x)=0$ 的两个实根，且 $x_1 < x_2$.

(1) 当 $x_1=a$ 或 $x_2=a$ 时，则 $x=a$ 不是 $f(x)$ 的极值点，此时不合题意.

(2) 当 $x_1 \ne a$ 且 $x_2 \ne a$ 时，由于 $x=a$ 是 $f(x)$ 的极大值点，故 $x_1 < a < x_2$.

即 $g(a) < 0$ 即 $a^2+(3-a+b)a+2b-ab-a < 0$

所以 $b < -a$ 所以 $b$ 的取值范围是 $(-\infty, -a)$. (步骤2)

(II) 由 (I) 可知，假设存在 $b$ 及 $x_4$ 满足题意，则

(1) 当 $x_2-a = a-x_1$ 时，则 $x_4 = 2x_2-a$ 或 $x_4 = 2x_1-a$,

于是 $2a = x_1+x_2 = a-b-3$. 即 $b = -a-3$.

此时 $x_4 = 2x_2-a = a-b-3+\sqrt{(a+b-1)^2+8}-a = a+2\sqrt{6}$

或 $x_4 = 2x_2-a = a-b-3-\sqrt{(a+b-1)^2+8}-a = a-2\sqrt{6}$ (步骤3)

(2) 当 $x_2-a = a-x_1$ 时，则 $x_2-a = 2(a-x_1)$ 或 $a-x_1 = 2(x_2-a)$,

\circled{1} 若 $x_2-a = 2(a-x_1)$, 则 $x_4 = \dfrac{a+x_2}{2}$,

于是 $3a = 2x_1+x_2 = \dfrac{3(a-b-3)-\sqrt{(a+b-1)^2+8}}{2}$,

即 $\sqrt{(a+b-1)^2+8} = -3(a+b+3)$, 于是 $a+b-1 = \dfrac{-9-\sqrt{13}}{2}$,

此时 $x_4 = \dfrac{a+x_2}{2} = \dfrac{2a+(a-b-3)-3(a+b+3)}{4} = -b-3 = a+\dfrac{1+\sqrt{13}}{2}$ (步骤4)

\circled{2} 若 $a-x_1 = 2(x_2-a)$, 则 $x_4 = \dfrac{a+x_2}{2}$

于是 $3a = 2x_2+x_1 = \dfrac{3(a-b-3)+\sqrt{(a+b-1)^2+8}}{2}$,

即 $\sqrt{(a-b-1)^2+8} = 3(a+b+3)$, 于是 $a+b-1 = \dfrac{-9+\sqrt{13}}{2}$,

此时 $x_4 = \dfrac{a+x}{2} = \dfrac{2a+(a-b-3)-3(a+b+3)}{4} = -b-3 = a+\dfrac{1-\sqrt{13}}{2}$ (步骤5)

综上所述，存在 $b$ 满足题意，

当 $b = -a-3$ 时， $x_4 = a \pm 2\sqrt{6}$,

当 $b = -a-\dfrac{7+\sqrt{13}}{2}$ 时， $x_4 = a+\dfrac{1+\sqrt{13}}{2}$,

当 $b = -a-\dfrac{7-\sqrt{13}}{2}$ 时， $x_4 = a+\dfrac{1-\sqrt{13}}{2}$. (步骤6)
\end{solutions}